\subsubsection{Urval}
\paragraph{}
Ett antal inklusions- och exklusionskriterier togs fram för att nå ett relevant resultat med litteraturstudien.
Artiklarna skulle vara publicerade i vetenskapliga tidskrifter som var peer-reviewed:ade och ha en kvalitativ ansats.


Söksträng: ('driver advisory system' AND railway) samt logiska operatorer.

\paragraph{}
Urvalsprocessen innehöll fyra steg.
Första steget innebar att utesluta verk baserat på titel.
Här sorterades eventuella dubbletter bort.
För andra steget lästes abstract och för tredje steget lästes utöver abstract sammanfattning/conclusion (motsv.).
Kvarvarande artiklar lästes i sin helhet.
För samtliga steg prövades artiklarna mot inklusions- och exklusionskriterierna för eventuell uteslutning.
Urvalsprocess en redovisas i tabell ~\ref{tab:Urvalsprocess}.
Utöver kvarvarande artiklar som återstod efter litteraturstudiens urvalsprocess ingick interna dokument från Trafikverket.
Sammanlagt inkluderades n verk i litteraturstudien.

\begin{table}
    \centering
    \begin{tabular}{|l|c|c|}
        \hline
        Urvalssteg & Databas & Antal träffar \\
        \hline
        1. Titel & IEEE & 408 \\
        2. Abstract & IEEE & 70 \\
        3. Slutsats & IEEE & 19 \\
        4. Hela artikeln & IEEE & 7 \\
        1. Titel & Google scholar & 408 \\
        2. Abstract & Google scholar & 70 \\
        3. Slutsats & Google scholar & 19 \\
        4. Hela artikeln & Google scholar & 7 \\
        \hline
    \end{tabular}
    \caption{Urvalsprocess}
    \label{tab:Urvalsprocess}
\end{table}



